\documentclass[11pt,a4paper]{article}
\usepackage[utf8]{inputenc}
\usepackage[T1]{fontenc}
\usepackage[margin=2.5cm]{geometry}
\usepackage{amsmath,amssymb,amsthm}
\usepackage{booktabs,array}
\usepackage{xcolor}
\usepackage{tcolorbox}
\tcbuselibrary{skins,breakable}
\usepackage{tikz}
\usetikzlibrary{arrows.meta,positioning,shapes.geometric,calc,fit,backgrounds}
\usepackage{hyperref}
\hypersetup{colorlinks=true,linkcolor=blue!60!black,urlcolor=blue!60!black,
            citecolor=blue!60!black}
\setlength{\emergencystretch}{2em}
\usepackage{enumitem}
\usepackage[numbers,sort&compress]{natbib}

\definecolor{boxblue}{RGB}{220,235,255}
\definecolor{boxorange}{RGB}{255,240,210}
\definecolor{boxgreen}{RGB}{220,245,220}
\definecolor{boxpurple}{RGB}{240,225,255}

\newcommand{\term}[1]{\textbf{#1}}
\newcommand{\lean}[1]{%
  \penalty-100
  \begingroup\def\_{\textunderscore\allowbreak}\texttt{#1}\endgroup}
\newcommand{\CatAL}{\ensuremath{\mathbf{A}_{\rm Lean}}}
\newcommand{\CatAD}{\textbf{A/D}}
\newcommand{\CatA}{\textbf{A}}
\newcommand{\CatB}{\textbf{B}}
\newcommand{\ZS}{\mathbb{Z}_7}
\newcommand{\PHIMDL}{$\Phi_{\rm MDL}$}
\newcommand{\Ngen}{N_{\rm gen}}
\newcommand{\cH}{c_H}
\newcommand{\sigW}{\sin^2\theta_W}

\newtheorem{theorem}{Theorem}[section]

\newcommand{\TutorialDOI}{10.5281/zenodo.20669144}
\date{2026}

\title{\textbf{The Forces:}\\[6pt]
\large How SU(3)$\times$SU(2)$\times$U(1) Emerges from the Polynomial\\[8pt]
\normalsize
\href{https://novaspivack.com/research}{novaspivack.com/research}\quad
\href{https://doi.org/10.5281/zenodo.20560550}{P48 complete theory monograph}\\[4pt]
\small\href{https://doi.org/\TutorialDOI}{doi.org/\TutorialDOI}
  \textit{(registration pending)}\\[4pt]
\normalsize\textit{UGP Physics Tutorial Series}}
\author{Nova Spivack}

\begin{document}
\setcounter{tocdepth}{2}
\maketitle
\tableofcontents
\newpage

\begin{tcolorbox}[colback=blue!6,colframe=blue!40!black,
  title=\textbf{UGP Physics Tutorial Series}]
\textbf{Prerequisites (read first):}
  \href{https://doi.org/10.5281/zenodo.20657898}{\emph{How the Standard Model Quantum Numbers Emerge}}
  $\cdot$
  \href{https://doi.org/10.5281/zenodo.20657923}{\emph{The Weinberg Angle from Orbit Arithmetic}}
  $\cdot$
  \href{https://doi.org/10.5281/zenodo.20657902}{\emph{The Strong CP Problem Solved by Structure}}\\[4pt]
\textbf{Next tutorials:}
  \href{https://doi.org/10.5281/zenodo.20657913}{\emph{From the Arithmetic Substrate to Particle Masses}}\\[4pt]
\textbf{See also:}
  \href{https://doi.org/10.5281/zenodo.20657889}{\emph{The GTE Polynomial: A Step-by-Step Cheat Sheet}}
  $\cdot$
  \href{https://doi.org/10.5281/zenodo.20657872}{\emph{The Levels of the Theory}}
\end{tcolorbox}

\begin{tcolorbox}[colback=boxgreen,colframe=green!50!black,
  title=\textbf{Claim-Strength Taxonomy}]
Throughout this tutorial, results are labeled with a confidence grade:
\begin{itemize}[itemsep=2pt]
  \item \textbf{$\mathbf{A}_{\rm Lean}$ (CatAL)} --- Machine-certified in Lean~4 with zero \texttt{sorry} and zero custom axioms.  Independently verifiable by anyone with Lean installed.
  \item \textbf{A/D (CatAD)} --- Analytically derived with full derivation, computationally confirmed.  Not yet machine-certified.
  \item \textbf{A (CatA)} --- Analytically derived; derivation complete but not independently machine-certified.
  \item \textbf{A/D (conditional)} --- Derived under a named, stated premise; the premise is itself an open problem or a CatB assumption.
  \item \textbf{B (CatB)} --- A stated working hypothesis or structural premise; not yet derived.
\end{itemize}
\end{tcolorbox}

%─────────────────────────────────────────────────────────────────────────────
\section{The Three Forces of the Standard Model}
\label{sec:forces-intro}
%─────────────────────────────────────────────────────────────────────────────

\begin{tcolorbox}[colback=boxblue,colframe=blue!40!black,title=The central claim]
In the Standard Model, the gauge group
$\mathrm{SU}(3) \times \mathrm{SU}(2) \times \mathrm{U}(1)$ is an assumption.
You write it down and then derive the forces from it.
In the GTE framework, the gauge group is a theorem.
You start from the 19-bit polynomial and derive which gauge group it implies ---
without putting any group in by hand.
\end{tcolorbox}

The three forces of the Standard Model (SM) are:

\begin{itemize}[itemsep=4pt]
  \item \textbf{Electromagnetism} (EM): carried by the photon; described by
        $\mathrm{U}(1)_{\rm em}$.
        It governs atomic structure, optics, and electric circuits.

  \item \textbf{The weak nuclear force}: carried by the $W^{\pm}$ and $Z^0$
        bosons; described by $\mathrm{SU}(2)_L \times \mathrm{U}(1)_Y$
        (the electroweak group, unified with EM in the Glashow--Weinberg--Salam
        model).
        It governs radioactive decay and neutrino interactions.
        The $L$ subscript signals \term{left-handedness}: the weak force
        acts only on left-handed particles.

  \item \textbf{The strong nuclear force}: carried by the gluons;
        described by $\mathrm{SU}(3)_c$ (the color group $c$).
        It binds quarks into protons and neutrons and holds nuclei together.
\end{itemize}

In the SM, the gauge group $\mathrm{SU}(3) \times \mathrm{SU}(2) \times \mathrm{U}(1)$
is postulated.
The theory is consistent with this choice, and consistent with many other choices
too.
There is no SM argument for why this particular group structure is the physical one.

The \term{Generative Triple Evolution (GTE)} framework derives it.
The derivation starts from the polynomial $p(L,C,R) = C + R - CR - LCR \pmod{7}$
and shows that its symmetry group structure uniquely implies
$\mathrm{SU}(3) \times \mathrm{SU}(2) \times \mathrm{U}(1)$ at the continuum level.

\subsection{What a gauge group actually does}

A \term{gauge group} encodes a local symmetry of a physical theory.
``Local'' means the symmetry can be applied differently at each point in spacetime.

The photon exists because electromagnetism has a local $\mathrm{U}(1)$ symmetry:
you can rotate the phase of a charged particle's wave function by an arbitrary
angle at each spacetime point, and the physics is unchanged.
To maintain this symmetry in the equations of motion, you must introduce a
compensating field --- the photon.

Similarly, the eight gluons of QCD exist because of the local $\mathrm{SU}(3)$
symmetry of the quark color charge.
The $W^{\pm}$ and $Z^0$ bosons exist because of the electroweak $\mathrm{SU}(2)_L
\times \mathrm{U}(1)_Y$ symmetry.

In each case, the gauge group determines:
(a) how many force-carrying particles exist,
(b) how they interact with matter, and
(c) how they interact with each other.

In the SM, (a)--(c) follow from the postulated group.
In GTE, the group itself is derived.

%─────────────────────────────────────────────────────────────────────────────
\section{The GTE Approach: $\mathbb{Z}_7$ Winding Conservation}
\label{sec:z7}
%─────────────────────────────────────────────────────────────────────────────

\begin{tcolorbox}[colback=boxgreen,colframe=green!50!black,
  title=The key structural fact]
At every Standard Model interaction vertex, the total $\mathbb{Z}_7$
winding number of the incoming particles equals the total $\mathbb{Z}_7$
winding number of the outgoing particles (mod 7).
This is a theorem, not a postulate.
Electric charge conservation, color conservation, and weak isospin
conservation all follow from it.
\end{tcolorbox}

\subsection{$\mathbb{Z}_7$ winding numbers}

The particles in the GTE framework are topological defects (kinks) in the
continuum field $\Phi_{\rm MDL}$ --- the physical field whose algebraic
certificate is the 19-bit polynomial.
Each particle is classified by its \term{winding number} in $\mathbb{Z}_7$
--- the group of integers mod 7, with elements $\{0, 1, 2, 3, 4, 5, 6\}$.

The winding number counts how many times the field winds around the
$\mathbb{Z}_7$ vacuum circle as you move through a particle.
It is conserved topologically --- you cannot remove a winding without
creating another.

The winding numbers of the SM particles, derived from the polynomial's
orbit structure (\CatAL)~\cite{SpivackGTECompleteFramework}, are:

\begin{table}[h]
\centering
\caption{$\mathbb{Z}_7$ winding numbers and SM force sectors.
All assignments machine-certified (\CatAL, \lean{winding\_class\_sm\_assignment}).}
\label{tab:winding}
\small
\begin{tabular}{cll}
\toprule
Winding $w$ & Particle sector & SM role \\
\midrule
0 & Vacuum, photon, neutrinos & Neutral / massless gauge sector \\
1 & Dark-mirror sector (PSC-forbidden) & No SM particle \\
2 & Quark (up-type: $u,c,t$) & Up-quark current \\
3 & $W^+$ boson, positron & Weak charged current \\
4 & Charged lepton ($e^-,\mu^-,\tau^-$), $W^-$ & Electromagnetic (QED) / weak charged current \\
5 & Dark-mirror sector (PSC-forbidden) & No SM particle \\
6 & Quark (down-type: $d,s,b$) & Charged weak, CKM mixing \\
\bottomrule
\end{tabular}
\end{table}

\subsection{Winding conservation at all 33 SM vertices}

In QED, the conservation of electric charge at every interaction vertex
is imposed as a requirement of the $\mathrm{U}(1)$ gauge symmetry.
In GTE, there is a prior conservation law: $\mathbb{Z}_7$ winding conservation
at every SM vertex is a theorem from the $\mathbb{Z}_7$ symmetry of the
$\Phi_{\rm MDL}$ Lagrangian.

\textbf{Worked example: the charged-current weak vertex.}
The $W^+$ boson connects an up-type quark ($w_u = 2$) to a down-type quark ($w_d = 6$):
\[
  u \to d + W^+.
\]
The $W^+$ carries winding $w_{W^+} = w_u - w_d \equiv 2 - 6 \equiv -4 \equiv 3 \pmod{7}$.
Check conservation: the incoming winding is $w_u = 2$; the outgoing winding is
$w_d + w_{W^+} = 6 + 3 = 9 \equiv 2 \pmod{7}$. \checkmark

\textbf{Worked example: electromagnetic vertex.}
An electron ($w = 4$) emits a photon ($w = 0$) and remains an electron ($w = 4$):
\[
  e^- \to e^- + \gamma.
\]
Incoming winding: 4. Outgoing winding: $4 + 0 = 4$. \checkmark

All 33 SM vertices have been verified by this calculation (\CatAL,
\lean{gte\_winding\_sm\_vertex\_conserved\_full})~\cite{SpivackGTECompleteFramework}.
Electric charge conservation is the projection of winding conservation
onto the degree-1 sector of the polynomial.

%─────────────────────────────────────────────────────────────────────────────
\section{The Frobenius Group $F_{21}$: The Symmetry Skeleton}
\label{sec:f21}
%─────────────────────────────────────────────────────────────────────────────

\subsection{From the polynomial to $F_{21}$}

The GTE polynomial $p(L,C,R) = C + R - CR - LCR \pmod{7}$ acts on the
finite field $\mathrm{GF}(7)$.
The symmetry group of this polynomial --- the group of transformations
that leave $p$ invariant --- is the \term{Frobenius group}
$F_{21} = \mathbb{Z}_7 \rtimes \mathbb{Z}_3$.

This is a finite group of order 21 (21 elements).
It is not a continuous (Lie) group.
The continuous SM gauge group $\mathrm{SU}(3) \times \mathrm{SU}(2) \times \mathrm{U}(1)$
operates at the continuum level.

The relationship between $F_{21}$ and the SM gauge group is:
$F_{21}$ is the discrete arithmetic skeleton; the continuous SM gauge group
is the continuum redundancy structure that the Algebraic Lifting Theorem
(\CatAL) forces to be present in the continuum field $\Phi_{\rm MDL}$,
given the symmetry structure of $F_{21}$.

The machine certification of $F_{21}$ as the symmetry group of the GTE polynomial
is \lean{frobenius\_f21\_is\_z7\_sdp\_z3} (\CatAL)~\cite{SpivackGTECompleteFramework}.

\subsection{The subgroup structure gives the forces}

$F_{21} = \mathbb{Z}_7 \rtimes \mathbb{Z}_3$ has the following key subgroup structure:
\begin{itemize}[itemsep=2pt]
  \item $\mathbb{Z}_7$ (order 7): the cyclic group of order 7.
        Its non-trivial action gives the color sector --- 7 distinct winding
        values, of which 3 correspond to color charge.
  \item $\mathbb{Z}_3$ (order 3): the cyclic group of order 3.
        Its action generates the 3-fold structure corresponding to color $\mathrm{SU}(3)$.
  \item The semidirect product structure $\rtimes$: encodes the interaction
        between the $\mathbb{Z}_7$ and $\mathbb{Z}_3$ sectors, giving the
        full gauge structure.
\end{itemize}

The MDL-minimality proof establishes that $\mathbb{Z}_7 \times \mathbb{Z}_3$
(or equivalently $F_{21}$) is the \emph{unique} structure of this form
with minimum description length:
\lean{mdl\_total\_z7z3\_strictly\_beats\_z5z3} (\CatAL)~\cite{SpivackGTECompleteFramework}.
Every competitor ($\mathbb{Z}_5 \times \mathbb{Z}_3$, $\mathbb{Z}_9 \times \mathbb{Z}_3$,
etc.) is either eliminated by no PSC-admissible kink orbits, or by higher MDL cost.

%─────────────────────────────────────────────────────────────────────────────
\section{The Strong Force: SU(3) from $F_{21}$}
\label{sec:strong}
%─────────────────────────────────────────────────────────────────────────────

\subsection{Color from the three ground states}

The GTE polynomial $p(L,C,R)$ over $\mathrm{GF}(7)$ has a set of ground states ---
values that the field can stably occupy.
The three non-trivial ground states are $\{0, 1, 5\}$ in $\mathrm{GF}(7)$.
These generate the group $\mathbb{Z}_3$ (a cyclic group of order 3: $\{0, 1, 2\} \pmod{3}$
in additive notation, or $\{1, 5, 25 \equiv 4\} \pmod{7}$ in multiplicative notation).

Three ground states, three colors.
The $\mathbb{Z}_3$ generated by the ground states is the discrete color charge.
At the continuum level, the MDL-minimal Lie group implementing $\mathbb{Z}_3$
color conservation is $\mathrm{SU}(3)$: the group of special unitary $3 \times 3$
matrices.
$\mathrm{SU}(3)$ contains $\mathbb{Z}_3$ as a subgroup (the center of $\mathrm{SU}(3)$)
and is the unique simple Lie group with a $\mathbb{Z}_3$ center at rank~2.

\subsection{Color confinement from MDL}

Why can we never observe a free quark?
In QCD, this \term{color confinement} is an experimental fact for which
no analytic proof from SM first principles exists (it is a Millennium Prize Problem).
In GTE, confinement follows from a description-length argument~\cite{SpivackGTECompleteFramework}.

A free colored quark in an asymptotic state requires $\log_2 3 + \log_2 15 = \log_2 45$
bits to describe (bits for the color choice plus bits for the winding value).
A color-neutral hadron requires only $\log_2 5$ bits (only 5 PSC-admissible winding values).
The description-length gap is:
\[
  \Delta K = \log_2 45 - \log_2 5 = \log_2 9 \approx 3.17~\text{bits} > 0.
\]
Because $\Delta K > 0$, a free colored quark is MDL-non-minimal.
PSC forbids MDL-non-minimal asymptotic states.
Therefore, no free colored quark exists.
Lean cert: \lean{color\_confinement\_k\_extra\_pos} (\CatAL, zero sorry)~\cite{SpivackGTECompleteFramework}.

\subsection{The QCD $\beta$-function coefficient $b_0 = 7$}

\term{Asymptotic freedom} is the property that the strong coupling $\alpha_s$ becomes
weaker at shorter distances (higher energies).
It follows from the QCD $\beta$-function being negative:
$\beta(\alpha_s) < 0 \iff b_0 > 0$.

In GTE, the one-loop $\beta$-function coefficient is:
\begin{equation}
  b_0 = |\mathbb{Z}_7| = 7.
  \label{eq:b0}
\end{equation}
This is a theorem from the Casimir relation of $F_{21}$
(\lean{b0\_eq\_z7\_order}, \CatAL)~\cite{SpivackGTECompleteFramework}.
Since $b_0 = 7 > 0$, the GTE substrate automatically gives an asymptotically free
color sector.
The QCD coupling is derived as $\alpha_s(M_Z) = 0.11822$ ($+0.24\sigma$
from PDG~2024, \CatAL).

%─────────────────────────────────────────────────────────────────────────────
\section{The Electroweak Force: SU(2)$_L \times$ U(1)$_Y$ from Chirality}
\label{sec:ew}
%─────────────────────────────────────────────────────────────────────────────

\subsection{The V$-$A structure: left-handedness from the CA rule pair}

The weak force has a distinctive property: it acts only on \term{left-handed}
particles.
This \term{V$-$A} (vector minus axial) structure is one of the SM's most striking
features.
In the SM, left-handedness is imposed by hand: the $\mathrm{SU}(2)_L$ group
acts only on the left-handed components of the fermion fields.

In GTE, left-handedness is structural.
The three-tape Chiral Minkowski Cellular Automaton (CMCA) has two outer layers:
Rule~110 (right-chiral) and Rule~124 (left-chiral).
These are not the same rule.
Rule~110 is related to Rule~124 by left-right mirror reflection.
The physical vacuum --- the unique $k = 0$ vacuum selected by the MDL principle
--- breaks the left-right symmetry: it selects one of the two chiral pair members
as the physical evolution rule for the outer tape.

This selection is the origin of parity violation.
The physical SM universe is left-handed because the MDL-selected vacuum
is left-handed.
Lean cert: \lean{sm\_orbit\_is\_left\_handed} (\CatAL)~\cite{SpivackGTECompleteFramework}.

\subsection{The Weinberg angle: $\sin^2\theta_W = 3/13$}

The \term{Weinberg angle} $\theta_W$ (also called the electroweak mixing angle)
is the angle that mixes the $\mathrm{SU}(2)_L$ and $\mathrm{U}(1)_Y$ gauge
bosons to produce the physical photon and $Z^0$.
Its value $\sin^2\theta_W \approx 0.231$ is a free parameter in the SM.

In GTE, it is the ratio of the generation count to the Higgs-sector orbit count~\cite{SpivackWeinbergAngle}:
\begin{equation}
  \sin^2\theta_W = \frac{\Ngen}{\cH} = \frac{3}{13} = 0.23077.
  \label{eq:weinberg}
\end{equation}
Here $\Ngen = 3$ is the number of SM generations, and $\cH = 13$ is the
Higgs color count: the number of distinct $\mathbb{Z}_7$ orbit sectors
accessible to the Higgs doublet under the $F_{21}$ action.

\textbf{Why $c_H = 13$?}
The Higgs doublet lives in the $F_{21}$ representation with $c_H$ accessible
channels.
The number 13 is not arbitrary: it is the next prime after the 12 electroweak
sector channels, and it is forced by the three-tape architecture.
The derivation is in P45~\cite{SpivackWeinbergAngle} and P48~\cite{SpivackGTECompleteFramework}.

\textbf{Comparison to experiment.}
The tree-level result $3/13 = 0.23077$ is a pure arithmetic result.
With two-loop threshold corrections (all inputs from GTE, no PDG parameters):
\[
  \sin^2\theta_W\big|_{\rm 2\text{-}loop} = 0.23129.
\]
PDG~2024 (MS-bar at $M_Z$): $\sin^2\theta_W = 0.23129 \pm 0.00004$.
Deviation: $+0.038\sigma$ (\CatAL/\CatAD)~\cite{SpivackGTECompleteFramework}.

\begin{tcolorbox}[colback=boxorange,colframe=orange!60!black,
  title={Result: Weinberg Angle (\CatAL)}]
Tree-level: $\sin^2\theta_W = \Ngen/\cH = 3/13 = 0.23077$\\[2pt]
Two-loop: $\sin^2\theta_W = 0.23129$ (\CatAD)\\[2pt]
PDG~2024: $0.23129 \pm 0.00004$. Deviation: $\mathbf{+0.038\sigma}$.\\[2pt]
Lean: \lean{weinberg\_angle\_closure}, \lean{weinberg\_angle\_derivation}
\end{tcolorbox}

\subsection{The $W$ and $Z$ boson masses}

Once the Higgs VEV $v_{\rm PSC} = 246.16$~GeV and the Weinberg angle are
derived, the gauge boson masses follow from standard electroweak relations:
\begin{align}
  M_W &= 80.364~\text{GeV} \quad (\CatAL;\;\lean{m\_W\_pdg\_interval}),
  \label{eq:mw}\\[4pt]
  M_Z &= M_W \cdot \sqrt{\frac{\cH}{\cH - \Ngen}}
       = M_W \cdot \sqrt{\frac{13}{10}}
       = 91.6~\text{GeV} \quad (\CatAL;\;\lean{m\_Z\_gte\_from\_weinberg}).
  \label{eq:mz}
\end{align}
PDG~2024 values: $M_W = 80.369 \pm 0.013$~GeV; $M_Z = 91.188 \pm 0.002$~GeV.
The GTE values are within $0.4\sigma$ and within $0.5\sigma$ respectively.

\subsection{Why the $Z$ is heavier than the $W$: the $c$-staircase}

The three electroweak bosons occupy a unit-step arithmetic staircase in the
$F_{21}$ branch-capacity parameter~\cite{SpivackGTECompleteFramework}:
\[
  c_{W^\pm} = 11, \quad c_Z = 12, \quad c_H = 13.
\]
Each step down from the Higgs scalar boundary $c_H = 13$ records the loss of
one longitudinal degree of freedom, encoding the Goldstone mechanism in
orbit arithmetic.
The photon (massless) sits at $w = 0$, the $\mathbb{Z}_7$ zero sector.
The $W$ is lighter than the $Z$ because it loses one more degree of freedom
($c_W < c_Z$).
No continuous parameter chooses these numbers; they are forced by the
integer structure of the $F_{21}$ orbits.
\mbox{Lean: \lean{ew\_c\_staircase} (\CatAL,\ zero sorry)}~\cite{SpivackGTECompleteFramework}.

%─────────────────────────────────────────────────────────────────────────────
\section{The Strong CP Problem: $\theta_{\rm QCD} = 0$ by Structure}
\label{sec:strongcp}
%─────────────────────────────────────────────────────────────────────────────

The SM Lagrangian admits a CP-violating term proportional to $\theta_{\rm QCD}$.
This term would cause the neutron electric dipole moment to be non-zero.
Experiment constrains $|\theta_{\rm QCD}| < 10^{-10}$ --- essentially exactly zero.
Why?

In GTE, $\theta_{\rm QCD} = 0$ exactly, with three independent proofs
from the structure of $F_{21}$~\cite{SpivackGTECompleteFramework}.

\textbf{Proof 1.} The $\theta$ term in the effective Lagrangian is proportional to
$\mathrm{Tr}[G_{\mu\nu} \tilde{G}^{\mu\nu}]$.
In $F_{21}$, this is a topological density.
The $\mathbb{Z}_7$ shift symmetry is anomaly-free at all loop orders
(\lean{z7\_shift\_measure\_preserving}, \CatAL), which forces the
$\theta$ term to vanish identically.

\textbf{Proof 2.} The $F_{21}$ group action on the gluon field configuration space
has no CP-violating fixed point.
The theorem \lean{f21\_theta\_term\_vanishes} (\CatAL) certifies this directly.

\textbf{Proof 3.} The three-tape architecture forces $\theta_{\rm QCD} = 0$ by
a discrete symmetry argument: the Rule~110/124 chiral pair has a discrete
CP symmetry that is preserved by the vacuum selection.

\begin{tcolorbox}[colback=boxorange,colframe=orange!60!black,
  title={Result: Strong CP Vanishes (\CatAL)}]
$\theta_{\rm QCD} = 0$ exactly. Three independent proofs.\\[2pt]
No axion required or predicted.\\[2pt]
Falsification: any measurement of a non-zero neutron electric dipole moment
at the $\theta_{\rm QCD} > 10^{-12}$ level would falsify the $F_{21}$
structure.\\[2pt]
Lean: \lean{f21\_theta\_term\_vanishes} (\CatAL, zero sorry)
\end{tcolorbox}

%─────────────────────────────────────────────────────────────────────────────
\section{The Fine Structure Constant: $\alpha^{-1} = 137$ at Tree Level}
\label{sec:alpha}
%─────────────────────────────────────────────────────────────────────────────

The \term{fine structure constant} $\alpha \approx 1/137.036$ governs the strength
of electromagnetic interactions.
In the SM, $\alpha$ is a free parameter fitted from experiment.
The value $\alpha^{-1} = 137.036$ at zero momentum transfer is measured to
extraordinary precision but not explained.

In GTE, the tree-level value is:
\begin{equation}
  \alpha^{-1} = 2^7 + 3^2 = 128 + 9 = 137.
  \label{eq:alpha}
\end{equation}
Lean cert: \lean{alpha\_em\_inverse\_structural\_identity} (\CatAL)~\cite{SpivackGTECompleteFramework}.

\textbf{Where do $2^7$ and $3^2$ come from?}
The factor $2^7 = 128$ counts the binary branching modes: the outer Rule~110 tape
has 7 binary generators, giving $2^7$ modes.
The factor $3^2 = 9$ counts the $\mathbb{Z}_3$ sector modes: the two-generation
residual from the three-generation orbit.
Both counts are forced by the polynomial's algebraic structure, not chosen to
reproduce $137$.

\textbf{Comparison to experiment.}
The tree-level GTE value is $\alpha^{-1} = 137$, while the measured value is
$\alpha^{-1} = 137.036$.
The fractional discrepancy is $0.026\%$, explained by QED running (radiative
corrections between zero momentum transfer and the physical scale).
The QED running brings the GTE result toward $137.036$ at \CatB~\cite{SpivackGTECompleteFramework}.

\textbf{A structural equivalence.}
The derivation can also be stated in terms of the GTE generation orbit.
The first-generation $b$-value $b_1 = 73$ satisfies
\[
  b_1 = 2^{\varphi(7)} + N_{\rm gen}^2,
\]
where $\varphi(7) = |\mathbb{F}_7^*| = 6$ is the order of the multiplicative group of
$\mathrm{GF}(7)$.  Since $\mathbb{F}_7^* \cong \mathbb{Z}_2 \times \mathbb{Z}_3$ and
$N_{\rm gen} = 3 = |\mathbb{Z}_3|$, the Sylow-3 subgroup of $\mathbb{F}_7^*$ has order $N_{\rm gen}$.
This gives $\varphi(7) = 2 \times N_{\rm gen}$, which is itself a structural identity.
The formula $\alpha^{-1} = 2^7 + N_{\rm gen}^2$ can therefore be rewritten as
$\alpha^{-1} = 2 \times b_1 - N_{\rm gen}^2 + N_{\rm gen}^2 = 2b_1$, or equivalently
checked directly: $2 \times 73 - 9 = 137$.
Both expressions --- the $|Z_7|$-mode count route $2^7 + N_{\rm gen}^2 = 137$ and the
orbit route via $b_1$ --- yield the same integer because $b_1 = 2^{\varphi(7)} + N_{\rm gen}^2$.
The two routes to $\alpha^{-1} = 137$ are thus algebraically equivalent under the
$\varphi(7) = 2N_{\rm gen}$ identity.
Lean cert: \lean{b1\_totient7\_ngen\_identity} (\CatAL)~\cite{SpivackGTECompleteFramework}.

%─────────────────────────────────────────────────────────────────────────────
\section{Electric Charge Quantization}
\label{sec:charge}
%─────────────────────────────────────────────────────────────────────────────

Electric charges in the SM come in multiples of $e/3$
(quark charges: $\pm 1/3$ and $\pm 2/3$; lepton charges: $0$ and $\pm 1$).
The SM requires charge quantization as a separate constraint.

In GTE, charge quantization is automatic.
The electric charge of any SM particle follows from the degree-1 projection
of the polynomial:
\begin{equation}
  Q = \frac{p(0, w, 0)}{3} = \frac{w_{\rm centered}}{3},
  \label{eq:charge}
\end{equation}
where $w_{\rm centered}$ is the centered representative of $w \in \mathbb{Z}_7$
in the range $[-3, 3]$.
The denominator 3 is not arbitrary: it equals the number of colors $N_c = 3$,
which is itself derived from the $\mathbb{Z}_3$ structure of $F_{21}$.
Charge is quantized in units of $e/3$ because there are exactly three colors.
A universe with two colors would have charge quantization in units of $e/2$;
a universe with four colors in units of $e/4$.
The GTE structure uniquely selects three colors, uniquely fixing the charge quantum.

\textbf{Worked examples:}
\begin{itemize}[itemsep=2pt]
  \item $w = 2$: $Q = 2/3$ (up quark). \checkmark
  \item $w = 4 \equiv -3$: $Q = -1$ (charged lepton). \checkmark
  \item $w = 3$: $Q = 1$ ($W^+$ boson). \checkmark
  \item $w = 0$: $Q = 0$ (neutrino, photon). \checkmark
  \item $w = 6 \equiv -1$: $Q = -1/3$ (down quark). \checkmark
\end{itemize}
Lean cert: \lean{gte\_charge\_formula} (\CatAL, zero sorry)~\cite{SpivackGTECompleteFramework}.
The denominator 3 in the charge formula comes from the color count $N_c = 3$,
which is itself derived from the polynomial's $\mathbb{Z}_3$ structure.

%─────────────────────────────────────────────────────────────────────────────
\section{Machine Certifications: The Complete Force Inventory}
\label{sec:inventory}
%─────────────────────────────────────────────────────────────────────────────

\begin{tcolorbox}[colback=boxpurple,colframe=blue!40!black,
  title={GTE force-sector machine certifications (\CatAL, zero sorry)},breakable]
\begin{itemize}[itemsep=2pt]
  \item \lean{gte\_winding\_sm\_vertex\_conserved\_full} ---
        $\mathbb{Z}_7$ conservation at all SM vertices
  \item \lean{gte\_charge\_formula} ---
        Charge quantization in multiples of $e/3$
  \item \lean{weinberg\_angle\_closure} ---
        $\sin^2\theta_W = N_{\rm gen}/c_H$ at tree level
  \item \lean{m\_W\_pdg\_interval} ---
        $M_W = 80.364$~GeV in PDG interval
  \item \lean{ew\_c\_staircase} ---
        Electroweak $c$-parameter staircase ($c_W = 11$, $c_Z = 12$, $c_H = 13$)
  \item \lean{color\_confinement\_k\_extra\_pos} ---
        Color confinement ($\Delta K = \log_2 9 > 0$)
  \item \lean{certified\_mass\_gap\_positive} ---
        QCD mass gap from $F_{21}$ orbit arithmetic
  \item \lean{b0\_eq\_z7\_order} ---
        QCD $\beta$-function $b_0 = |\mathbb{Z}_7| = 7$
  \item \lean{f21\_theta\_term\_vanishes} ---
        $\theta_{\rm QCD} = 0$ (strong CP)
  \item \lean{alpha\_em\_inverse\_structural\_identity} ---
        $\alpha^{-1} = 2^7 + 3^2 = 137$
  \item \lean{sm\_orbit\_is\_left\_handed} ---
        Left-handedness from vacuum selection
  \item \lean{gravity\_em\_degree\_split} ---
        Gravity (degree 3) vs.\ EM (degree 1) from the polynomial
\end{itemize}
\end{tcolorbox}

\subsection{How a finite group produces continuous gauge symmetry}

A natural question: $F_{21}$ is a finite group of 21 elements.
How does it produce a continuous gauge theory with infinitely many local
symmetry transformations?

The answer is that $F_{21}$ and the SM gauge group operate at different levels.
At Level~1 (the algebraic certificate, the CMCA), $\mathbb{Z}_7$ winding
conservation is an exact discrete arithmetic law.
At Level~2 (the continuum $\Phi_{\rm MDL}$ field), the Algebraic Lifting
Theorem (\CatAL) requires that the same conservation laws hold continuously.
Local continuous redundancy of those discrete constraints is the gauge principle.
The MDL-minimal Lie group implementing the $\mathbb{Z}_7 \times \mathbb{Z}_3$
discrete structure is $\mathrm{SU}(3) \times \mathrm{SU}(2) \times \mathrm{U}(1)$.

The discrete group is the arithmetic skeleton.
The continuous gauge group is the continuum redundancy structure.

%─────────────────────────────────────────────────────────────────────────────
\section*{Further Reading}
%─────────────────────────────────────────────────────────────────────────────
\begin{tcolorbox}[colback=orange!8,colframe=orange!50!black,
  title=\textbf{Primary Sources}]
The results in this tutorial are derived in the following papers,
all available open-access on Zenodo and at
\href{https://novaspivack.com/research}{novaspivack.com/research}:
\begin{itemize}[itemsep=2pt]
  \item \textbf{P48 --- The Complete GTE Framework} (main synthesis monograph):\\
        \href{https://doi.org/10.5281/zenodo.20560550}{doi.org/10.5281/zenodo.20560550}
  \item \textbf{P45 --- The Three-Tape Chiral Minkowski Cellular Automaton}
        (Weinberg angle, chirality, force structure):\\
        \href{https://doi.org/10.5281/zenodo.20465805}{doi.org/10.5281/zenodo.20465805}
  \item \textbf{P08 --- From UGP to GTE: Prime-Locked Universes}
        (MDL selection, $F_{21}$ group theory):\\
        \href{https://doi.org/10.5281/zenodo.20171002}{doi.org/10.5281/zenodo.20171002}
  \item \textbf{P01 --- Standard Model Parameters from the UGP}
        (coupling constant derivations):\\
        \href{https://doi.org/10.5281/zenodo.20168787}{doi.org/10.5281/zenodo.20168787}
\end{itemize}
Lean~4 machine proofs: \href{https://github.com/novaspivack/ugp-lean}{github.com/novaspivack/ugp-lean}
\end{tcolorbox}

\bibliographystyle{unsrtnat}
\bibliography{../../papers/bib/Spivack_Papers_Bibliography}

\end{document}
